% ======================================================================
\subsection{PCP decomposition and the NEGF self-energy}
\label{sub:pcp_selfenergy}
% ======================================================================

We analyse whether the Permanent CP (PCP) decomposition of Luo et
al.~\cite{luo_tensor_2025} can yield a self-energy kernel that is both
correct and cheaper than dense evaluation.


% -------------------------------------------------------
\subsubsection{PCP ansatz}
\label{sssub:pcp_ansatz}
% -------------------------------------------------------

The PCP decomposition expresses the mass-weighted FC3 as
\begin{equation}\label{eq:pcp_ansatz}
  \Psi_{\alpha\beta\gamma}(0\kappa,\, l'\kappa',\, l''\kappa'')
  \;=\; \sum_{\xi=1}^{N_c} \frac{\lambda_\xi}{3!}
  \sum_{\sigma\in S_3}
    A_{\sigma_1}^\xi(0,\kappa,\alpha)\;
    A_{\sigma_2}^\xi(l',\kappa',\beta)\;
    A_{\sigma_3}^\xi(l'',\kappa'',\gamma),
\end{equation}
where $A_i^\xi(\vc{R},\kappa,\alpha)$ are three real-valued mode
functions (``legs'') for each rank~$\xi$, and $S_3$ denotes the six
permutations of the three legs.  Each leg is a \emph{vector} in the
combined (cell, atom, Cartesian) index space.

Fourier-transforming each mode,
\begin{equation}\label{eq:pcp_ft}
  \hat{f}_i^\xi(\vq)_a
  = \sum_{\vc{R},\kappa}
    \frac{e^{-i\vq\cdot\vc{R}}}{\sqrt{M_\kappa}}\;
    A_i^\xi(\vc{R},\kappa,\alpha),
  \qquad a = 3\kappa + \alpha,
\end{equation}
gives $\hat{f}_i^\xi(\vq) \in \CC^M$ where $M = 3n_\mathrm{at}$ is
the number of primitive DOFs.  The Fourier-transformed FC3 vertex is
\begin{equation}\label{eq:pcp_vertex}
  \hat\Phi_{acd}(\vq_1,\vq_2)
  = \sum_{\xi} \frac{\lambda_\xi}{3!}
    \sum_{\sigma\in S_3}
    \hat{f}_{\sigma_1}^\xi(\vq_\mathrm{ext})_a\;
    \hat{f}_{\sigma_2}^\xi(\vq_1)_c\;
    \hat{f}_{\sigma_3}^\xi(\vq_2)_d\,.
\end{equation}


% -------------------------------------------------------
\subsubsection{Substitution into the self-energy}
\label{sssub:pcp_substitution}
% -------------------------------------------------------

The phonon--phonon self-energy (\cref{eq:sigma_phph}) has the form
\begin{equation}\label{eq:sigma_structure}
  \Sigma^{\gtrless}_{ab}(\omega;\vq)
  = \frac{i\hbar}{2N}\sum_{\vq'}
    \sum_{c,d,e,f}
    \hat\Phi_{acd}(\vq',\vq{-}\vq')\;
    K_{cdfe}(\omega;\vq',\vq{-}\vq')\;
    \hat\Phi^*_{bef}(\vq{-}\vq',\vq'),
\end{equation}
where
$K_{cdfe}(\omega;\vq_1,\vq_2) = \int \frac{d\omega'}{2\pi}\,
G^{\gtrless}_{cf}(\omega';\vq_1)\,G^{\gtrless}_{de}(\omega{-}\omega';\vq_2)$
is the Green's function frequency convolution.  The index pairing
in~$K$ is crucial:
\begin{itemize}
  \item $G(\vq')$ carries indices $(c, f)$,
  \item $G(\vq - \vq')$ carries indices $(d, e)$.
\end{itemize}

Substituting the PCP vertex \cref{eq:pcp_vertex} for the left
factor and its conjugate for the right factor:
\begin{align}\label{eq:pcp_sigma_full}
  \Sigma^{\gtrless}_{ab}(\omega;\vq)
  &= \frac{i\hbar}{2N}\sum_{\vq'}
    \sum_{\xi,\xi'} \frac{\lambda_\xi \lambda_{\xi'}}{36}
    \sum_{\sigma,\sigma'\in S_3}
    \hat{f}_{\sigma_1}^\xi(\vq)_a\;
    \hat{f}_{\sigma'_1}^{\xi'\,*}(\vq)_b
    \notag\\
  &\quad\times
    \sum_{c,d,e,f}
    \hat{f}_{\sigma_2}^\xi(\vq')_c\;
    \hat{f}_{\sigma_3}^\xi(\vq{-}\vq')_d\;
    K_{cdfe}\;
    \hat{f}_{\sigma'_2}^{\xi'\,*}(\vq{-}\vq')_e\;
    \hat{f}_{\sigma'_3}^{\xi'\,*}(\vq')_f\,.
\end{align}
The external legs $\hat{f}_{\sigma_1}^\xi(\vq)_a$ and
$\hat{f}_{\sigma'_1}^{\xi'\,*}(\vq)_b$ depend only on $\vq$ (not
$\vq'$) and can be pulled out of the $\vq'$-sum.


% -------------------------------------------------------
\subsubsection{Factorization of the internal contraction}
\label{sssub:pcp_factorization}
% -------------------------------------------------------

The internal contraction over $(c,d,e,f)$ factors because the
index pairing in $K$ aligns with the PCP mode structure.
Expanding $K_{cdfe} = \int G_{cf}(\omega';\vq')\,G_{de}(\omega{-}\omega';\vq{-}\vq')\,d\omega'/2\pi$:
\begin{align}
  &\sum_{c,d,e,f}
    f_{\sigma_2,c}\,f_{\sigma_3,d}\,
    G_{cf}(\omega';\vq')\,G_{de}(\omega{-}\omega';\vq{-}\vq')\,
    f^*_{\sigma'_2,e}\,f^*_{\sigma'_3,f}
  \notag\\
  &\quad=
    \underbrace{
      \Bigl[\sum_c f_{\sigma_2,c}\;\sum_f G_{cf}(\omega';\vq')\;
      f^*_{\sigma'_3,f}\Bigr]
    }_{g_{\sigma_2,\sigma'_3}^{\xi\xi'}(\omega';\vq')}
    \;\times\;
    \underbrace{
      \Bigl[\sum_d f_{\sigma_3,d}\;\sum_e G_{de}(\omega{-}\omega';\vq{-}\vq')\;
      f^*_{\sigma'_2,e}\Bigr]
    }_{g_{\sigma_3,\sigma'_2}^{\xi\xi'}(\omega{-}\omega';\vq{-}\vq')},
  \label{eq:factorization}
\end{align}
where we define the scalar Green's function projections:
\begin{equation}\label{eq:scalar_gf}
  g_{ij}^{\xi\xi'}(\omega;\vq)
  \;=\; \hat{f}_i^{\xi}(\vq)^\dagger\,
    G^{\gtrless}(\omega;\vq)\,
    \hat{f}_j^{\xi'}(\vq).
\end{equation}
The factorization \cref{eq:factorization} is \emph{exact} because
indices $(c, f)$ both belong to $G(\vq')$ and indices $(d, e)$
both belong to $G(\vq - \vq')$---the two Green's functions
never share an internal index.

The full PCP self-energy is therefore:
\begin{equation}\label{eq:pcp_sigma_scalar}
  \boxed{
  \Sigma^{\gtrless}_{ab}(\omega;\vq)
  = \frac{i\hbar}{2N}\sum_{\vq'}
    \sum_{\xi,\xi'} \frac{\lambda_\xi \lambda_{\xi'}}{36}
    \sum_{\sigma,\sigma'}
    \hat{f}_{\sigma_1}^\xi(\vq)_a\;
    \underbrace{
      \mathrm{IFFT}_\omega\Bigl\{
        \hat{g}_{\sigma_2,\sigma'_3}^{\xi\xi'}(\vq')\;\cdot\;
        \hat{g}_{\sigma_3,\sigma'_2}^{\xi\xi'}(\vq{-}\vq')
      \Bigr\}
    }_{\text{scalar convolution}}
    \;\hat{f}_{\sigma'_1}^{\xi'\,*}(\vq)_b\,.
  }
\end{equation}

\paragraph{Permutation sum simplification.}
The code groups the double sum over $S_3 \times S_3$ by
$(\sigma_1, \sigma'_1)$ indices, forming
\begin{equation}
  \mathcal{S}[\xi,\xi',s_1,s_1',\omega]
  = \sum_{\substack{(s_1,s_2,s_3)\in S_3 \\ (s_1',s_2',s_3')\in S_3}}
    g_{s_2,s_2'}(\vq')\;\cdot\; g_{s_3,s_3'}(\vq{-}\vq'),
\end{equation}
which uses index pairing $(s_2, s_2')$ and $(s_3, s_3')$ rather
than the derived pairing $(s_2, s_3')$ and $(s_3, s_2')$.  These
are \textbf{equivalent}: summing over all permutations of
$\{0,1,2\}\setminus\{s_1\}$ and
$\{0,1,2\}\setminus\{s_1'\}$ generates all four combinations of
the remaining two elements on both sides, so both orderings
produce the same sum.  Numerically verified to machine precision
(Test~4, relative difference $\sim 10^{-17}$).


% -------------------------------------------------------
\subsubsection{Comparison with the separable (SVD) factorization}
\label{sssub:pcp_vs_svd_structure}
% -------------------------------------------------------

Both the PCP and separable decompositions lead to \emph{exact}
factorizations of the self-energy internal contraction, but for
different structural reasons.

\paragraph{SVD.}
$\hat\Phi_{acd}(\vq_1,\vq_2) = \sum_r F_{r,ac}(\vq_1)\,h_{r,d}(\vq_2)$
has a \emph{matrix} left factor $F_r \in \CC^{M\times M}$ and a
\emph{vector} right factor $h_r \in \CC^M$.  The Green's function
$G(\vq')$ with indices $(c,f)$ is contracted:
$F_{r,ac}\,G_{cf}\,h^*_{s,f} = [F_r\,G\,h_s^*]_a \in \CC^M$
(a vector), and $G(\vq-\vq')$ with indices $(d,e)$:
$h_{r,d}\,G_{de}\,F^*_{s,be} = [h_r^T G F_s^\dagger]_b \in \CC^M$.
The frequency convolution is between two $M$-vectors, yielding
an $M \times M$ outer product---correct by construction.

\paragraph{PCP.}
$\hat\Phi_{acd}$ has \emph{vector} factors on all three indices.
The Green's function projections yield \emph{scalars}
$g = f^\dagger G f \in \CC$.  The frequency convolution is between
two scalars, and the result is expanded back to $M \times M$ via
the external mode vectors.  This is also correct because the
index pairing in $K_{cdfe}$ guarantees that all four internal
indices are covered by mode vectors at the correct $\vq$-points.


% -------------------------------------------------------
\subsubsection{The Fourier transform problem}
\label{sssub:ft_problem}
% -------------------------------------------------------

While the self-energy factorization is exact, the \emph{vertex}
computation via PCP mode FTs is not.  The PCP real-space ansatz
\cref{eq:pcp_ansatz} involves a sum over cell index~$l$ with
modular arithmetic:
\begin{equation}
  \Psi(0, l', l'')
  = \sum_\xi \frac{\lambda_\xi}{6} \sum_\sigma \sum_{l=0}^{N-1}
    A_{\sigma_1}\bigl((-l) \bmod N\bigr)\;
    A_{\sigma_2}\bigl((l'-l) \bmod N\bigr)\;
    A_{\sigma_3}\bigl((l''-l) \bmod N\bigr).
\end{equation}
Fourier-transforming $\Psi$ over $(l', l'')$ and factorizing into
individual mode FTs requires
\begin{equation}\label{eq:phase_identity}
  e^{-i\vq\cdot\vc{R}_{(m+l)\bmod N}}
  \;=\; e^{-i\vq\cdot\vc{R}_m}\;e^{-i\vq\cdot\vc{R}_l},
\end{equation}
which holds if and only if
$\vq \cdot N_i\vc{a}_i = 2\pi n_i$ for each lattice
direction~$i$, i.e., $\vq$ is \emph{commensurate} with the
supercell.  When $(m + l) \bmod N$ wraps around a boundary,
$\vc{R}_{(m+l)\bmod N} = \vc{R}_m + \vc{R}_l - N_i\vc{a}_i$
for some direction~$i$, introducing a spurious phase
$e^{i\vq\cdot N_i\vc{a}_i} \neq 1$.

For a $4\times 4$ $q$-mesh on a $2\times 2\times 2$ supercell,
12 of 16 $q$-points are non-commensurate.  The vertex error
at these points reaches 50--110\% (\cref{tab:vertex_errors}).

\begin{table}[htbp]
  \centering
  \caption{PCP vertex errors: factorized FT vs.\ correct dense FT
    (from PCP-reconstructed FC3). Commensurate $q$-points have
    $\vq \cdot N_i\vc{a}_i \in 2\pi\ZZ$ for all~$i$.}
  \label{tab:vertex_errors}
  \begin{tabular}{llrl}
    \toprule
    $\vq'$ & $\vq-\vq'$ & $\|\hat\Phi_\mathrm{PCP} - \hat\Phi_\mathrm{dense}\|/\|\hat\Phi_\mathrm{dense}\|$ & Status \\
    \midrule
    $(0, \tfrac{1}{2})$       & $(0, \tfrac{1}{2})$       & $6.5\times10^{-16}$ & Commensurate \\
    $(\tfrac{1}{2}, 0)$       & $(\tfrac{1}{2}, \tfrac{1}{2})$ & $7.7\times10^{-16}$ & Commensurate \\
    \midrule
    $(0, \tfrac{1}{4})$       & $(0, \tfrac{3}{4})$       & $5.1\times10^{-1}$ & Non-commensurate \\
    $(\tfrac{1}{4}, \tfrac{1}{4})$ & $(\tfrac{3}{4}, \tfrac{1}{2})$ & $1.1\times10^{0}$ & Non-commensurate \\
    \bottomrule
  \end{tabular}
\end{table}


% -------------------------------------------------------
\subsubsection{Numerical verification: factorized PCP is exact}
\label{sssub:pcp_numerical}
% -------------------------------------------------------

To isolate the FT issue from the factorization question, we
test at commensurate $q$-points where the mode FTs are exact.

\paragraph{Test setup.}
Si $2\times2\times2$ supercell, PCP fit with $N_c = 24$
(rel.\ error $3.9\times10^{-4}$).  Random anti-Hermitian
Green's functions $G^<, G^>$ on a $2\times2$ $q$-mesh (all 4
points commensurate) and 31 frequency points.

\paragraph{Self-energy comparison.}
\begin{center}
\begin{tabular}{lc}
  \toprule
  Component & $\|\Sigma_\mathrm{PCP} - \Sigma_\mathrm{dense}\| / \|\Sigma_\mathrm{dense}\|$ \\
  \midrule
  $\Sigma^<$ & $8.7\times10^{-16}$ \\
  $\Sigma^>$ & $9.3\times10^{-16}$ \\
  $\Sigma^R$ & $9.1\times10^{-16}$ \\
  \bottomrule
\end{tabular}
\end{center}
At every $q$-point individually, the error is $\sim 10^{-16}$
(machine precision).  \textbf{The factorized PCP self-energy
is mathematically exact}---the error is entirely in the Fourier
transform, not the contraction structure.

On a $4\times4$ $q$-mesh (12 non-commensurate points), the
self-energy error rises to $\sim 94\%$.

\paragraph{Why it works.}
The self-energy contraction $\Sigma_{ab} = \Phi^L_{acd}\,K_{cdfe}\,\Phi^R_{bef}$
involves $K_{cdfe} = \mathrm{IFFT}\{G_{cf}(\vq')\,G_{de}(\vq-\vq')\}$.
The PCP mode vectors project $G(\vq')$ via indices $(c,f)$ and
$G(\vq-\vq')$ via indices $(d,e)$---the two Green's functions
never share an internal index.  Therefore the scalar projections
$g = f^\dagger G f$ are applied to independent tensor factors,
and the sums commute exactly.


% -------------------------------------------------------
\subsubsection{Cost analysis}
\label{sssub:pcp_cost}
% -------------------------------------------------------

Since the factorized PCP self-energy is exact (given correct vertices),
the cost comparison with dense and separable kernels is meaningful.

\paragraph{Per-pair cost structure.}
\begin{enumerate}
  \item \textbf{Scalar projections} (precomputed per $\vq$):
    $g_{ij}^{\xi\xi'}(\omega;\vq) = \hat{f}_i^{\xi\dagger}\,G(\omega;\vq)\,\hat{f}_j^{\xi'}$.
    Cost: $3N_c$ matrix-vector products $(M^2)$ + $9N_c^2$ inner
    products $(M)$, over $N_\omega$ frequencies.

  \item \textbf{Permutation sums and FFT convolution}: For each
    $(\vq,\vq')$ pair, form the $S_3 \times S_3$ permutation sums
    and convolve via FFT.
    Cost: $36 N_c^2 N_\tau + 9 N_c^2 N_\tau\log N_\tau$.

  \item \textbf{Accumulation}: Expand to matrix self-energy via
    $\Sigma_{ab} \mathrel{+}= f_a\,[C \cdot f_b^*]$.
    Cost: $9 N_\omega (N_c^2 M + N_c M^2)$.
\end{enumerate}

Total per-pair:
\begin{equation}\label{eq:pcp_cost}
  \cW_\mathrm{PCP}^\mathrm{pair}
  = \underbrace{9 N_c^2\,N_\tau \log N_\tau}_{\text{FFT}}
  + \underbrace{9 N_\omega\,(N_c^2 M + N_c M^2)}_{\text{accumulation}}.
\end{equation}
Compare:
\begin{align}
  \cW_\mathrm{dense}^\mathrm{pair}
  &= M^4\,N_\tau\log N_\tau + N_\omega\,M^5,
  \label{eq:dense_cost_recap} \\
  \cW_\mathrm{sep}^\mathrm{pair}
  &= R^2\,N_\tau\log N_\tau + R^2\,N_\omega\,M^2.
  \label{eq:sep_cost_recap}
\end{align}

\paragraph{PCP vs.\ dense.}
PCP is cheaper when $9N_c^2 < M^4$, i.e., $N_c < M^2/3$.
For Si ($M = 6$): $N_c < 12$ gives a speedup in the FFT step.
At $N_c = 4$: FFT cost $144$ vs.\ $1296$ ($9\times$ cheaper);
accumulation $9(96 + 144) = 2160$ vs.\ $7776$ ($3.6\times$).

\paragraph{PCP vs.\ separable.}
At equal compression ($N_c = R$), the PCP FFT cost carries a
factor of 9 from the permutation structure.  However, the PCP
accumulation is $N_c^2 M + N_c M^2$ vs.\ SVD's $R^2 M^2$, which
is cheaper when $N_c < M$ (since $N_c M + M^2 < R M^2$ for
$N_c < M$).  For Si at $R = N_c = 6$: PCP accumulation
$= 9(216 + 216) = 3888$ vs.\ SVD $= 36 \cdot 36 = 1296$.  The
factor of 9 overhead makes PCP slower at $M = 6$.

For large primitive cells ($M \gg N_c$), the PCP accumulation
term $9 N_c M^2$ matches SVD's $R^2 M^2$ when $9 N_c \approx R^2$,
i.e., at $R \approx 3\sqrt{N_c}$.


% -------------------------------------------------------
\subsubsection{Resolving the Fourier transform problem}
\label{sssub:ft_resolution}
% -------------------------------------------------------

Three approaches to obtain correct PCP vertices at arbitrary
$\vq$-points:

\paragraph{1.\ Commensurate $q$-mesh.}
Use a $q$-mesh whose grid spacing divides the supercell, so all
$q$-points satisfy \cref{eq:phase_identity}.  For a $2\times2\times2$
supercell, this limits the mesh to $2\times2$.  Coarser than
desirable for converged results but trivially correct.

\paragraph{2.\ Real-space reconstruction.}
Reconstruct FC3 from PCP modes, build $M_\mathrm{stacked}$ and
gathering matrices $T(\vq)$, compute vertices via
$\hat\Phi_a(\vq_1,\vq_2) = T(\vq_1)\,M_a\,T(\vq_2)^T$.  This is
the current implementation.  Correct at arbitrary $\vq$ but uses
the dense FT machinery, losing the $N_c^2$ cost advantage.

\paragraph{3.\ Direct vertex evaluation.}
For each $(\vq_1, \vq_2)$, evaluate the PCP vertex by summing
over cell index~$l$ \emph{without} factorizing the phases:
\begin{equation}\label{eq:direct_vertex}
  \hat\Phi_{acd}(\vq_1,\vq_2)
  = \sum_\xi \frac{\lambda_\xi}{6} \sum_\sigma \sum_{l=0}^{N-1}
    A_{\sigma_1}((-l)\bmod N)_a\;
    \Bigl[\sum_m A_{\sigma_2}(m)_c\,
      e^{-i\vq_1\cdot\vc{R}_{(m+l)\bmod N}}\Bigr]\;
    \Bigl[\sum_n A_{\sigma_3}(n)_d\,
      e^{-i\vq_2\cdot\vc{R}_{(n+l)\bmod N}}\Bigr].
\end{equation}
The bracketed ``shifted FTs'' have cost $\cO(N_\mathrm{cell} \cdot M)$
each.  For $N_c$ modes, 3 legs, $N_\mathrm{cell}$ shifts, and
6 permutations: $6 N_c \cdot N_\mathrm{cell} \cdot 2 N_\mathrm{cell} M$
$= 12 N_c N_\mathrm{cell}^2 M$.  This must be done for each
of the $N_q^2$ pairs $(\vq', \vq-\vq')$.

For Si ($N_c = 8$, $N_\mathrm{cell} = 8$, $M = 6$, $N_q = 16$):
vertex cost per pair $\approx 12 \cdot 8 \cdot 64 \cdot 6 = 36864$,
vs.\ dense vertex $M \cdot (d_t^2 + d_t M) = 6 \cdot 2592 = 15552$.
Direct vertex is $\sim 2.4\times$ more expensive than the
$M_\mathrm{stacked}$ approach per pair.

However, the PCP vertex can be \emph{precomputed} once and
cached for all SCBA iterations.  If the SCBA loop runs for
$K$ iterations, the vertex precomputation cost is amortized
and the per-iteration cost reduces to the scalar convolution
\cref{eq:pcp_cost}---potentially much cheaper than dense.

\paragraph{4.\ Supercell-commensurate interpolation.}
Compute the factorized vertex exactly at the
$N_\mathrm{cell}$ commensurate $q$-points and interpolate to the
finer mesh.  The vertex is smooth in $\vq$-space (it is the DFT
of a finite-range real-space object), so trigonometric
interpolation is exact up to the supercell cutoff range.  This
is equivalent to zero-padding the real-space FC3 and DFT'ing,
i.e., the reconstruction approach but potentially more memory-efficient
for large systems where storing the full FC3 is prohibitive.


% -------------------------------------------------------
\subsubsection{Summary}
\label{sssub:pcp_summary}
% -------------------------------------------------------

\begin{enumerate}
  \item The factorized PCP self-energy
    (\cref{eq:pcp_sigma_scalar}) is \textbf{mathematically exact}.
    The scalar projection of $G$ onto PCP mode vectors commutes
    with the frequency convolution because the two Green's functions
    in $K_{cdfe}$ carry disjoint index pairs $(c,f)$ and $(d,e)$
    that each pair with mode vectors at a single $\vq$-point.

  \item The sole source of error is the \textbf{Fourier transform}:
    the PCP mode FTs factorize only at $q$-points commensurate with
    the supercell.  At non-commensurate points, the modular
    cell-index arithmetic produces phase errors of $\cO(1)$.

  \item With correct vertices (via direct evaluation or real-space
    reconstruction), the PCP kernel at $N_c = 4$ would be
    $\sim 9\times$ cheaper than dense in the FFT step and
    $\sim 3.6\times$ cheaper in accumulation (for Si, $M = 6$).

  \item The current implementation uses real-space reconstruction
    (approach~2), which is correct but loses the cost advantage.
    Implementing the direct vertex evaluation (approach~3) with
    SCBA-amortized caching would enable the factorized PCP kernel
    at arbitrary $q$-points with the correct cost scaling.

  \item For the NEGF self-energy, the relevant comparison is with
    the separable (SVD) kernel, which also gives exact results and
    has cost $R^2 M^2$ vs.\ PCP's $9 N_c^2 M$.  The factor of 9
    from $S_3$ permutations makes PCP slower at small $M$, but for
    large primitive cells ($M \gg N_c$), the $N_c M$ accumulation
    term becomes favorable if $9 N_c < R^2$.
\end{enumerate}
